\subsection{Setup}
\label{sec:setup}

\paragraph{Platform.} Lifecycle and SDK runs use the stock Android~13 emulator image
(android-33 google\_apis arm64-v8a rev.~17, build TE1A.240213.009), emulator 37.2.13, on an 18-core Apple-silicon host. The stream-wrapper campaign uses two AVDs side by side; the enlarged matched campaign uses four read-only instances of each of those two configurations. HTTP/TLS cost and label-error tests use a separate Android~15 emulator on the same host; exact fingerprints and host load are archived.
AVD~A runs \vname{none}, \vname{cancel}, \vname{enqueue}, \vname{jobstart}, \vname{conv} and \vname{rd}. AVD~B runs
\vname{egress}, \vname{rd}, \vname{volatile} and \vname{taint}. Repeating \vname{rd} on both AVDs does not remove variant--AVD confounding: cross-AVD comparisons are descriptive. Matched channel variants run within each AVD. Before every trial each package is cleared with \texttt{pm clear}, granted location, opened once,
and placed in the \texttt{active} standby bucket. Without that last step JobScheduler never ran our jobs: a never-opened app
sits in the \texttt{NEVER} bucket.

\begin{table}[t]
\centering
\caption{\lab tenants. Every tenant has 12 tagged records; ``sum.'' tenants also derive a summary (declared scope in brackets).}
\label{tab:tenants}
\mtab\setlength{\tabcolsep}{2.5pt}
\begin{tabular}{@{}l>{\raggedright\arraybackslash}p{1.85cm}>{\raggedright\arraybackslash}p{1.75cm}>{\raggedright\arraybackslash}p{2.9cm}@{}}
\toprule
\mhead Tenant & Contract & Withdrawal & What it tests \\
\midrule
T1 & opt-out, discard & in-app & summary filed as \textit{diagnostics} (lineage) \\
T8 & opt-out, discard & in-app & same, summary correctly declared \\
T2 & deletion request & in-app + delete rows & summary copy in job extras \\
T5 & opt-out, discard & withdraw, re-grant & epoch vs.\ boolean (+4 new records) \\
T9 & opt-out, discard & in-app & raw-socket upload (uncovered path) \\
T10 & opt-out, discard & in-app & encrypted body (unattributable) \\
\midrule
T3 & opt-out, flush allowed & in-app & FP control: queued data may go \\
T4 & restriction & personalisation & FP control: backup must go \\
T6 & none & none & FP control: everything must go \\
T7 & OS permission & \texttt{pm revoke} (in \textsf{rev}) & FP control: prior fixes may sync \\
\bottomrule
\end{tabular}
\end{table}

\paragraph{Schedules.} Each trial sets up all ten tenants in every package, injects one schedule, waits until no job of
ours remains registered (polling \texttt{dumpsys jobscheduler}), and then reads each app's event log with \texttt{sqlite3}
as root, which does not start the app. The seven schedules are:
\begin{itemize}[leftmargin=*,nosep]
\item \textsf{pen}: withdraw while the jobs are pending.
\item \textsf{kill}: withdraw, then \texttt{kill -9}.
\item \textsf{fstop}: withdraw, force-stop, reopen the app 3\,s later.
\item \textsf{reb}: withdraw, then reboot.
\item \textsf{rev}: withdraw, then \texttt{pm revoke} location, which kills the app.
\item \textsf{net}: connectivity cut as the jobs start and restored 6\,s later; withdrawal 1.5--1.8\,s after job start.
\item \textsf{infl}: withdrawal at a seeded uniform offset in $[0,1.6]$\,s after the upload job starts.
\end{itemize}
In \textsf{net} and \textsf{infl}, jobs are started with \texttt{cmd jobscheduler run -f}, and the withdrawal is injected
inside the app relative to the job's own start, so its timing is precise. In \textsf{net}, each app gets its own 6\,s outage in
turn, during which only the trigger uploads of T8, T2, T5 and T4 are started. The outage is a guest firewall rule rejecting traffic to the
server. We first tried \texttt{svc wifi/data disable}, but it did not cut the apps' traffic reliably, and starting all apps' jobs
at once let JobScheduler's concurrency limit decide which variant saw the outage. The configuration trials using those two methods are excluded. Schedule parameters are derived from a seed per
(schedule, trial), so both AVDs replay identical schedules.

\paragraph{Oracle and metrics.} Ground truth comes from the sink. Tags are visible in plaintext payloads, and a summary's
content lists its source tags. The stream-wrapper device sequence orders withdrawals and pre-write authorization logs,
not actual byte acceptance. Baselines can therefore undercount forbidden handoffs in the
log-to-write gap. Missing sequence evidence is ambiguous; host receipt time alone cannot
reconstruct a socket-handoff order. We preserve the authorization-order estimand explicitly.
The channel oracle uses paired write-attempt/accepted-count events and the withdrawal
transaction order; interrupted or overlapping observations remain unknown. We report:
\begin{itemize}[leftmargin=*,nosep]
\item \emph{Violation rate}: the fraction of tenant-trials with at least one prohibited item, with trial-cluster uncertainty, keeping all tenants and paired variants together.
\item \emph{Prohibited items} per tenant-trial.
\item \emph{Legit deficit}: permitted items that a variant had not delivered when the observation window closed (\vname{none}'s deficit is the availability loss that enforcement does not cause).
\item \emph{Unattributable} submissions: post-withdrawal bodies the sink cannot read.
\item \emph{Outstanding} submissions: accepted before withdrawal, possibly received afterwards; stream-wrapper logs identify only authorization/response overlap.
\end{itemize}
A trial is excluded if the host slept during it (any gap between harness actions over 200\,s), if the emulator restarted,
or if any device log could not be read (\S\ref{sec:threats}). Of \nTrialsTotal{} logged trials, \nInvalid{} were excluded.
We analyze the first 50 valid trials per schedule and AVD (\nTrialsValid{} trials); the \nValidUnused{} further valid trials,
run as replacements while exclusions were still being counted, are kept in the artifact but not used.

Trial uncertainty uses 10,000 trial-cluster bootstrap resamples (seed 20261115), stratified by AVD and schedule. The 350 stream-wrapper AVD-A trials give a conventional-baseline authorization rate of 26.0\% (percentile interval [25.8\%, 26.3\%]). This is not a socket-handoff rate. A zero-event percentile interval is descriptive, not a zero-risk bound. All tenants and paired variants exposed to a trial remain one resampling unit.

